\documentclass[11pt]{article}
\usepackage[a4paper,margin=1in]{geometry}
\usepackage{booktabs}
\usepackage{amsmath}
\usepackage{amssymb}
\usepackage{hyperref}

\title{BONK Bullish L2 Orderbook Research: Current Conclusion}
\author{Finance Chain Research Workspace}
\date{2026-05-13}

\begin{document}
\maketitle

\section{Scope}

The current canonical run tag is
\[
  \texttt{20260513\_bullish\_l2\_basket\_price\_v1}.
\]
The data window is 2026-04-29 through 2026-05-12. The main BONK venues are
\[
  \texttt{BONK1MUSDC}, \qquad \texttt{BONK1MUSDT}.
\]
The analysis uses Bullish L2 state, path labels, cross-venue state, and Binance spot price context. It does not use options data, does not use incremental L2 reconstruction, and does not claim a trading rule or alpha.

\section{Main Finding}

The strongest current object is not a global predictive model. It is a conditional orderbook state:
\[
  \texttt{BONK1MUSDT, H4: depth\_high + rv\_low}
\]
and the cleaner extension
\[
  \texttt{BONK1MUSDT, H4: depth\_high + rv\_low + cv\_spread}.
\]

Interpretation:
\[
  \text{deep book} + \text{quiet BONK volatility} + \text{cross-venue state}
  \Rightarrow
  \text{better 4h path quality}.
\]
``Better path quality'' mainly means stronger residual-positive behavior and lower-first suppression. It does not mean that each selected minute is an immediate buy signal.

\section{What The Orderbook Is Telling Us}

The orderbook data is most useful as a state and path-risk layer:
\begin{enumerate}
  \item \textbf{State filter.} High displayed depth and low realized volatility identify a deep and quiet regime.
  \item \textbf{Path filter.} The useful label is first-passage behavior: whether the path hits the upper or lower barrier first.
  \item \textbf{Execution constraint.} Spread and displayed depth decide whether any gross edge can survive costs.
  \item \textbf{Cross-venue control.} USDC/USDT disagreement is useful as venue state, but not as a clean lead-lag signal.
\end{enumerate}

The correct research chain is therefore
\[
  \text{orderbook state}
  \rightarrow
  \text{path labels}
  \rightarrow
  \text{market/meme controls}
  \rightarrow
  \text{capacity and cost checks}.
\]
The less reliable chain is
\[
  \text{all orderbook features}
  \rightarrow
  \text{one global model}
  \rightarrow
  \text{trade signal}.
\]

\section{Key Evidence}

\begin{table}[h]
\centering
\begin{tabular}{lll}
\toprule
Check & Result & Read \\
\midrule
Exact label audit & 0.205\% strict mismatch & Labels do not explain H4 gate \\
Active H4 gate impact & 0 / 10 reads changed & H4 conclusion stable to exact labels \\
BONK1MUSDT coverage & 99.84\% minutes & L2 coverage is strong \\
Primary H4 gate affected rows & 0 / 1127 & Missing minutes do not explain gate \\
Cross-venue lead-lag & max 0.0047 Spearman & No clear USDC/USDT leader \\
Feature-family ablation & cross-venue ranks best & Useful as state/control \\
Combined orderbook model & underperforms & Too many L2 features overfit \\
\bottomrule
\end{tabular}
\caption{High-level robustness checks.}
\end{table}

\section{Regime Interpretation}

The V4 state taxonomy places the active gate mostly in a deep and quiet regime:
\[
  \texttt{deep\_quiet\_wide\_common\_off\_cv\_mixed}.
\]
For the active combo gate
\[
  \texttt{depth\_high + rv\_low + cv\_spread},
\]
about 82.0\% of selected rows fall into that state.

This means the gate is not simply a tight-spread liquidity state. It is deeper and quieter than average, but the spread can be wide and cross-venue state can be mixed. In plain terms, BONK looks most interesting when the book is deep enough to support movement, volatility is quiet enough to reduce lower-first path risk, and USDC/USDT venue state is not fully aligned.

\section{Path And Edge Read}

The current H4 state has enough gross path movement to deserve research follow-up. Previous gross-edge diagnostics showed that the H4 active gates can survive simple spread/fee stress on paper. The V4 capacity check refines this:
\[
  \text{top-book envelope} \approx 0,
  \qquad
  \text{5-level envelope} \approx 50 \text{ quote},
  \qquad
  \text{25-level stress envelope} \approx 500 \text{ quote}.
\]
So the phenomenon is not zero-capacity, but it is small. It should remain a research candidate, not an execution-ready strategy.

\section{Short-Horizon Read}

Short-horizon labels are feasible from existing data:
\[
  H \in \{5\text{m},15\text{m},30\text{m}\}, \qquad
  B \in \{20,30\}\text{ bps}.
\]
For 5m labels, 100 bps is too wide and mostly misses the microstructure question. The useful short-horizon work should focus on 20--30 bps barriers.

The strongest short-horizon diagnostics come from spread, depth, activity, and microprice state. However, adding all orderbook features to a context model does not consistently improve proper scores. This again supports selective state use instead of feature dumping.

\section{Current BONK Status}

The current BONK situation can be summarized as:
\[
\boxed{
\begin{aligned}
&\text{BONK has a plausible H4 orderbook-state phenomenon;}\\
&\text{the phenomenon is regime-local, not a universal model edge;}\\
&\text{it is not explained away by exact-label or missing-minute issues;}\\
&\text{it is partially cross-venue / market-state dependent;}\\
&\text{capacity is small and must be treated as a hard constraint.}
\end{aligned}
}
\]

The active candidate remains:
\[
  \texttt{BONK1MUSDT H4 depth\_high + rv\_low + cv\_spread}.
\]
The candidate should be classified as
\[
  \texttt{research\_state},
\]
not
\[
  \texttt{trade\_rule}.
\]

\section{Blockers}

Promotion is blocked by:
\begin{enumerate}
  \item Negative-control weakness: random phase and time-shift controls are still too strong.
  \item Short data window: only about two weeks of Bullish L2 data are available.
  \item Capacity: displayed size is small, especially at top of book.
  \item Model instability: gate-internal models and combined feature models do not consistently beat the manual state.
  \item No incremental L2 queue reconstruction: strict OFI and queue-level behavior are not yet tested.
\end{enumerate}

\section{Next Step}

The next useful engineering step is:
\[
  \text{canonical exact labels}
  +
  \text{canonical short-horizon labels}
  \rightarrow
  \text{rerun frozen V4 checks}.
\]
More L2 data collection is justified only if the frozen state survives:
\[
  \text{exact labels}
  \land
  \text{negative controls}
  \land
  \text{market/meme/SOL controls}
  \land
  \text{displayed-depth stress}.
\]

\end{document}
